%! AP Statistics — Unit 2, Lesson 2.1 — Warm-Up KEY
\documentclass[10pt]{article}
\usepackage{apstats-article}
\usepackage{apstats-key}

%=============================================================================
\begin{document}

\pageheader{Unit 2, Lesson 2.1}{Warm-Up --- Answer Key}
\begin{tabularx}{\linewidth}{X X X}
Name:\;\ans{Key} & Date:\; & Period:\;
\end{tabularx}

\vspace{0.3in}

{\large\bfseries\color{navy} Part 1 \quad Spiral Review}

\vspace{0.12in}

\begin{enumerate}

    \item Classify each variable.

    \vspace{8pt}
    \renewcommand{\arraystretch}{1.8}
    \begin{tabularx}{\linewidth}{@{} X c c @{}}
    \textbf{Variable} & \textbf{Categorical} & \textbf{Quantitative} \\
    \hline
    Preferred genre of music        & \ans{\checkmark} & \\
    Hours of sleep per night        & & \ans{\checkmark} \\
    Whether a student plays a sport & \ans{\checkmark} & \\
    GPA                             & & \ans{\checkmark} \\
    \end{tabularx}

    \vspace{0.18in}

    \item \ans{Accept any reasonable research question relating two variables from the table.
    Example: ``Do students who play sports get more or fewer hours of sleep per night than
    students who do not play sports?'' or ``Is there a relationship between GPA and hours
    of sleep per night?''}

\end{enumerate}

\vspace{0.25in}

{\large\bfseries\color{navy} Part 2 \quad Looking Ahead}

\vspace{0.12in}

\begin{tcolorbox}[colback=greenbg, colframe=greenacc, boxrule=0.8pt, arc=2mm,
    left=4mm, right=4mm, top=3mm, bottom=3mm]
    \small
    \begin{itemize}[leftmargin=1.3em, itemsep=8pt]
  \item An apparent relationship in data may be \textbf{random} --- not every pattern is real.
  \item Even a real relationship does not mean one variable \textbf{causes} the other.
  \item The variable types (categorical or quantitative) determine which analysis we use
        in Lessons 2.2--2.9.
    \end{itemize}
\end{tcolorbox}

\vspace{0.2in}

\begin{teachernote}
\textbf{Part 1 notes:} GPA and hours of sleep are quantitative; music genre and sport
participation are categorical. Students who mark GPA as categorical are likely thinking of
letter grades --- clarify that GPA is a number we can average, not a category.

\textbf{Part 2:} No scored answer; check that students wrote something (engagement check).
Common responses: ``What is a confounding variable?'' or ``What does `association' mean?''
--- both are great previews for today's lesson.
\end{teachernote}

\end{document}
